% Ngân hàng bài mở rộng: Hằng đẳng thức đáng nhớ (lớp 8), bài 032 đến 051.
% Quy ước:
%   - Mỗi bài: begin{problem}[id] ... end{problem}; lời giải lồng bên trong.
%   - id: duy nhất, không dấu, không đổi, không dùng lại. Dạng <lớp><phân môn>-<chủ đề>-<số>.
%   - Mỗi bước lời giải là 1 đoạn, các bước cách nhau 1 dòng trống; mỗi bước có lý do.
%   - Dòng cuối lời giải: Đáp số: $\boxed{...}$
%   - Cách giải khác: thêm begin{solution}[cách 2].
%   - Metadata bằng comment ngay sau begin{problem}: % nguon:, % lop:, % kc_du_kien:, % do_kho_du_kien:
\documentclass[12pt,a4paper]{article}
\usepackage[utf8]{vietnam}
\usepackage{amsmath,amssymb}
\newenvironment{problem}[1][]{\par\medskip\noindent\textbf{Bài #1.}\ }{\par}
\newenvironment{solution}[1][]{\par\noindent\textit{Lời giải #1.}\par}{\par}

\begin{document}

% ---------------------------------------------------------------- viết đa thức thành bình phương
\begin{problem}[8DS-HDT-032]
% nguon: dummy (tự soạn để mở rộng ngân hàng)
% lop: 8
% kc_du_kien: 8DS-HDT-01
% do_kho_du_kien: 1
Viết mỗi đa thức sau dưới dạng bình phương của một tổng hoặc một hiệu:

a) $x^2 + 8x + 16$;

b) $9y^2 - 12y + 4$;

c) $25a^2 + 30ab + 9b^2$.

\begin{solution}
a) Ta có $16 = 4^2$ và $8x = 2 \cdot x \cdot 4$, nên đa thức có dạng $A^2 + 2AB + B^2$ với $A = x$, $B = 4$. Theo hằng đẳng thức bình phương của một tổng: $x^2 + 8x + 16 = (x + 4)^2$.

b) Ta có $9y^2 = (3y)^2$, $4 = 2^2$ và $12y = 2 \cdot 3y \cdot 2$, nên đa thức có dạng $A^2 - 2AB + B^2$ với $A = 3y$, $B = 2$. Theo hằng đẳng thức bình phương của một hiệu: $9y^2 - 12y + 4 = (3y - 2)^2$.

c) Ta có $25a^2 = (5a)^2$, $9b^2 = (3b)^2$ và $30ab = 2 \cdot 5a \cdot 3b$, nên đa thức có dạng $A^2 + 2AB + B^2$ với $A = 5a$, $B = 3b$. Theo hằng đẳng thức bình phương của một tổng: $25a^2 + 30ab + 9b^2 = (5a + 3b)^2$.

Đáp số: a) $\boxed{(x + 4)^2}$; b) $\boxed{(3y - 2)^2}$; c) $\boxed{(5a + 3b)^2}$
\end{solution}
\end{problem}

% ---------------------------------------------------------------- khai triển lập phương
\begin{problem}[8DS-HDT-033]
% nguon: dummy (tự soạn để mở rộng ngân hàng)
% lop: 8
% kc_du_kien: 8DS-HDT-03
% do_kho_du_kien: 1
Khai triển các biểu thức sau:

a) $(x + 4)^3$;

b) $(2y - 3)^3$;

c) $(a + 2b)^3$.

\begin{solution}
a) Áp dụng hằng đẳng thức lập phương của một tổng $(A + B)^3 = A^3 + 3A^2B + 3AB^2 + B^3$ với $A = x$, $B = 4$: $(x + 4)^3 = x^3 + 3 \cdot x^2 \cdot 4 + 3 \cdot x \cdot 4^2 + 4^3 = x^3 + 12x^2 + 48x + 64$.

b) Áp dụng hằng đẳng thức lập phương của một hiệu $(A - B)^3 = A^3 - 3A^2B + 3AB^2 - B^3$ với $A = 2y$, $B = 3$: $(2y - 3)^3 = (2y)^3 - 3 \cdot (2y)^2 \cdot 3 + 3 \cdot 2y \cdot 3^2 - 3^3 = 8y^3 - 36y^2 + 54y - 27$.

c) Áp dụng hằng đẳng thức lập phương của một tổng với $A = a$, $B = 2b$: $(a + 2b)^3 = a^3 + 3 \cdot a^2 \cdot 2b + 3 \cdot a \cdot (2b)^2 + (2b)^3 = a^3 + 6a^2b + 12ab^2 + 8b^3$.

Đáp số: a) $\boxed{x^3 + 12x^2 + 48x + 64}$; b) $\boxed{8y^3 - 36y^2 + 54y - 27}$; c) $\boxed{a^3 + 6a^2b + 12ab^2 + 8b^3}$
\end{solution}
\end{problem}

% ---------------------------------------------------------------- tính nhanh, ba hằng đẳng thức
\begin{problem}[8DS-HDT-034]
% nguon: dummy (tự soạn để mở rộng ngân hàng)
% lop: 8
% kc_du_kien: 8DS-HDT-05, 8DS-HDT-02, 8DS-HDT-01, 8DS-HDT-03
% do_kho_du_kien: 2
Tính nhanh:

a) $75^2 - 25^2$;

b) $53^2 + 94 \cdot 53 + 47^2$;

c) $101^3 - 3 \cdot 101^2 + 3 \cdot 101 - 1$.

\begin{solution}
a) Áp dụng hằng đẳng thức hiệu hai bình phương $A^2 - B^2 = (A - B)(A + B)$ với $A = 75$, $B = 25$: $75^2 - 25^2 = (75 - 25)(75 + 25) = 50 \cdot 100 = 5\,000$.

b) Viết $94 = 2 \cdot 47$, khi đó biểu thức là $53^2 + 2 \cdot 53 \cdot 47 + 47^2$, có dạng $A^2 + 2AB + B^2$ với $A = 53$, $B = 47$. Theo hằng đẳng thức bình phương của một tổng: $(53 + 47)^2 = 100^2 = 10\,000$.

c) Biểu thức có dạng $A^3 - 3A^2B + 3AB^2 - B^3$ với $A = 101$, $B = 1$ (vì $3 \cdot 101^2 \cdot 1$ và $3 \cdot 101 \cdot 1^2$). Theo hằng đẳng thức lập phương của một hiệu: $(101 - 1)^3 = 100^3 = 1\,000\,000$.

Đáp số: a) $\boxed{5\,000}$; b) $\boxed{10\,000}$; c) $\boxed{1\,000\,000}$
\end{solution}
\end{problem}

% ---------------------------------------------------------------- rút gọn biểu thức
\begin{problem}[8DS-HDT-035]
% nguon: dummy (tự soạn để mở rộng ngân hàng)
% lop: 8
% kc_du_kien: 8DS-HDT-01, 8DS-HDT-02, 8DS-HDT-04
% do_kho_du_kien: 2
Rút gọn các biểu thức sau:

a) $(x + 3)^2 - (x - 3)^2$;

b) $(2a - b)^2 + (2a + b)^2 - 2(4a^2 - b^2)$;

c) $(x + 2)(x^2 - 2x + 4) - (x^3 - 1)$.

\begin{solution}
a) Áp dụng hằng đẳng thức bình phương của một tổng và của một hiệu: $(x + 3)^2 = x^2 + 6x + 9$ và $(x - 3)^2 = x^2 - 6x + 9$. Bỏ ngoặc (trước ngoặc có dấu trừ thì đổi dấu các hạng tử) rồi thu gọn: $x^2 + 6x + 9 - x^2 + 6x - 9 = 12x$.

b) Áp dụng các hằng đẳng thức bình phương: $(2a - b)^2 = 4a^2 - 4ab + b^2$ và $(2a + b)^2 = 4a^2 + 4ab + b^2$; nhân đơn thức với đa thức: $2(4a^2 - b^2) = 8a^2 - 2b^2$.

b) Thay vào và thu gọn các hạng tử đồng dạng: $4a^2 - 4ab + b^2 + 4a^2 + 4ab + b^2 - 8a^2 + 2b^2 = (4a^2 + 4a^2 - 8a^2) + (-4ab + 4ab) + (b^2 + b^2 + 2b^2) = 4b^2$.

c) Áp dụng hằng đẳng thức tổng hai lập phương $(A + B)(A^2 - AB + B^2) = A^3 + B^3$ với $A = x$, $B = 2$: $(x + 2)(x^2 - 2x + 4) = x^3 + 8$. Do đó biểu thức bằng $x^3 + 8 - x^3 + 1 = 9$.

Đáp số: a) $\boxed{12x}$; b) $\boxed{4b^2}$; c) $\boxed{9}$
\end{solution}
\end{problem}

% ---------------------------------------------------------------- bài toán thực tế: tính nhanh diện tích
\begin{problem}[8DS-HDT-036]
% nguon: dummy (tự soạn để mở rộng ngân hàng)
% lop: 8
% kc_du_kien: 8DS-HDT-05, 8DS-HDT-02, 8DS-HDT-01
% do_kho_du_kien: 2
Một thửa ruộng hình chữ nhật có chiều dài $105$ m và chiều rộng $95$ m. Bên cạnh có một sân phơi hình vuông cạnh $97$ m. Không dùng máy tính, hãy tính:

a) Diện tích thửa ruộng;

b) Diện tích sân phơi;

c) Diện tích nào lớn hơn và lớn hơn bao nhiêu mét vuông.

\begin{solution}
a) Viết $105 = 100 + 5$ và $95 = 100 - 5$, rồi áp dụng hằng đẳng thức hiệu hai bình phương: $105 \cdot 95 = (100 + 5)(100 - 5) = 100^2 - 5^2 = 10\,000 - 25 = 9\,975$ (m$^2$).

b) Viết $97 = 100 - 3$ và áp dụng hằng đẳng thức bình phương của một hiệu: $97^2 = (100 - 3)^2 = 100^2 - 2 \cdot 100 \cdot 3 + 3^2 = 10\,000 - 600 + 9 = 9\,409$ (m$^2$).

c) So sánh hai kết quả: $9\,975 > 9\,409$ nên thửa ruộng có diện tích lớn hơn sân phơi, và chênh lệch là $9\,975 - 9\,409 = 566$ (m$^2$).

Đáp số: a) $\boxed{9\,975 \text{ m}^2}$; b) $\boxed{9\,409 \text{ m}^2}$; c) $\boxed{\text{thửa ruộng lớn hơn } 566 \text{ m}^2}$
\end{solution}
\end{problem}

% ---------------------------------------------------------------- chứng minh đẳng thức bằng tổng, hiệu hai lập phương
\begin{problem}[8DS-HDT-037]
% nguon: dummy (tự soạn để mở rộng ngân hàng)
% lop: 8
% kc_du_kien: 8DS-HDT-04, 8DS-HDT-06
% do_kho_du_kien: 2
Chứng minh đẳng thức $(x + y)(x^2 - xy + y^2) + (x - y)(x^2 + xy + y^2) = 2x^3$ với mọi số $x$, $y$.

\begin{solution}
Áp dụng hằng đẳng thức tổng hai lập phương $(A + B)(A^2 - AB + B^2) = A^3 + B^3$ với $A = x$, $B = y$: $(x + y)(x^2 - xy + y^2) = x^3 + y^3$.

Áp dụng hằng đẳng thức hiệu hai lập phương $(A - B)(A^2 + AB + B^2) = A^3 - B^3$ với $A = x$, $B = y$: $(x - y)(x^2 + xy + y^2) = x^3 - y^3$.

Cộng hai kết quả và thu gọn: vế trái bằng $x^3 + y^3 + x^3 - y^3 = 2x^3$, đúng bằng vế phải (điều phải chứng minh).

Đáp số: $\boxed{(x + y)(x^2 - xy + y^2) + (x - y)(x^2 + xy + y^2) = 2x^3}$
\end{solution}
\end{problem}

% ---------------------------------------------------------------- tính tổng và hiệu khi biết tổng bình phương và tích
\begin{problem}[8DS-HDT-038]
% nguon: dummy (tự soạn để mở rộng ngân hàng)
% lop: 8
% kc_du_kien: 8DS-HDT-05, 8DS-HDT-01
% do_kho_du_kien: 2
Cho hai số dương $x$, $y$ thoả mãn $4x^2 + 9y^2 = 25$, $xy = 2$ và $2x > 3y$.

a) Tính $(2x + 3y)^2$ và $(2x - 3y)^2$.

b) Tính $2x + 3y$ và $2x - 3y$.

\begin{solution}
a) Áp dụng hằng đẳng thức bình phương của một tổng với $A = 2x$, $B = 3y$: $(2x + 3y)^2 = 4x^2 + 12xy + 9y^2 = (4x^2 + 9y^2) + 12xy = 25 + 12 \cdot 2 = 49$.

a) Áp dụng hằng đẳng thức bình phương của một hiệu: $(2x - 3y)^2 = 4x^2 - 12xy + 9y^2 = (4x^2 + 9y^2) - 12xy = 25 - 12 \cdot 2 = 1$.

b) Vì $x > 0$, $y > 0$ nên $2x + 3y > 0$; số dương có bình phương bằng $49$ là $7$, do đó $2x + 3y = 7$.

b) Vì $2x > 3y$ nên $2x - 3y > 0$; số dương có bình phương bằng $1$ là $1$, do đó $2x - 3y = 1$.

Đáp số: a) $\boxed{(2x + 3y)^2 = 49,\ (2x - 3y)^2 = 1}$; b) $\boxed{2x + 3y = 7,\ 2x - 3y = 1}$
\end{solution}
\end{problem}

% ---------------------------------------------------------------- bài toán thực tế: chu vi và diện tích hình chữ nhật
\begin{problem}[8DS-HDT-039]
% nguon: dummy (tự soạn để mở rộng ngân hàng)
% lop: 8
% kc_du_kien: 8DS-HDT-01, 8DS-HDT-05
% do_kho_du_kien: 3
Một mảnh vườn hình chữ nhật có chu vi $28$ m và diện tích $45$ m$^2$. Biết chiều dài lớn hơn chiều rộng.

a) Không tìm riêng từng kích thước, hãy tính hiệu giữa chiều dài và chiều rộng.

b) Tính chiều dài và chiều rộng của mảnh vườn.

\begin{solution}
a) Gọi chiều dài và chiều rộng lần lượt là $a$ (m) và $b$ (m), điều kiện $a > b > 0$. Chu vi bằng $28$ m nên $2(a + b) = 28$, suy ra $a + b = 14$; diện tích bằng $45$ m$^2$ nên $ab = 45$.

a) Áp dụng hai hằng đẳng thức bình phương: $(a + b)^2 = a^2 + 2ab + b^2$ và $(a - b)^2 = a^2 - 2ab + b^2$. Trừ vế theo vế: $(a + b)^2 - (a - b)^2 = 4ab$, suy ra $(a - b)^2 = (a + b)^2 - 4ab$.

a) Thay số: $(a - b)^2 = 14^2 - 4 \cdot 45 = 196 - 180 = 16$. Vì $a > b$ nên $a - b > 0$; số dương có bình phương bằng $16$ là $4$, do đó $a - b = 4$.

b) Từ $a + b = 14$ và $a - b = 4$, cộng vế theo vế được $2a = 18$, $a = 9$; suy ra $b = 14 - 9 = 5$.

b) Đối chiếu điều kiện: $9 > 5 > 0$ (thoả mãn). Thử lại: chu vi $2(9 + 5) = 28$ m và diện tích $9 \cdot 5 = 45$ m$^2$ (đúng).

Đáp số: a) $\boxed{4 \text{ m}}$; b) $\boxed{9 \text{ m và } 5 \text{ m}}$
\end{solution}
\end{problem}

% ---------------------------------------------------------------- giá trị nhỏ nhất, giá trị lớn nhất của tam thức bậc hai
\begin{problem}[8DS-HDT-040]
% nguon: dummy (tự soạn để mở rộng ngân hàng)
% lop: 8
% kc_du_kien: 8DS-HDT-06, 8DS-HDT-01
% do_kho_du_kien: 3
a) Tìm giá trị nhỏ nhất của biểu thức $A = 4x^2 + 4x + 5$.

b) Tìm giá trị lớn nhất của biểu thức $B = 7 + 6x - 9x^2$.

\begin{solution}
a) Tách $5 = 1 + 4$ để làm xuất hiện bình phương của một tổng: $A = (4x^2 + 4x + 1) + 4 = (2x)^2 + 2 \cdot 2x \cdot 1 + 1^2 + 4 = (2x + 1)^2 + 4$.

a) Vì $(2x + 1)^2 \geq 0$ với mọi $x$ nên $A \geq 4$. Dấu bằng xảy ra khi và chỉ khi $2x + 1 = 0$, tức là $x = -\dfrac{1}{2}$. Thử lại: $A = 4 \cdot \dfrac{1}{4} - 2 + 5 = 4$ (đúng).

b) Đặt dấu trừ ra ngoài các hạng tử chứa $x$ rồi thêm và bớt $1$: $B = -(9x^2 - 6x) + 7 = -(9x^2 - 6x + 1) + 1 + 7 = -(3x - 1)^2 + 8$.

b) Vì $(3x - 1)^2 \geq 0$ nên $-(3x - 1)^2 \leq 0$, do đó $B \leq 8$. Dấu bằng xảy ra khi và chỉ khi $3x - 1 = 0$, tức là $x = \dfrac{1}{3}$. Thử lại: $B = 7 + 2 - 1 = 8$ (đúng).

Đáp số: a) $\boxed{\min A = 4 \text{ khi } x = -\dfrac{1}{2}}$; b) $\boxed{\max B = 8 \text{ khi } x = \dfrac{1}{3}}$
\end{solution}
\end{problem}

% ---------------------------------------------------------------- tính a^3 - b^3, ab, a^2 + b^2 khi biết a - b và a^2 + ab + b^2
\begin{problem}[8DS-HDT-041]
% nguon: dummy (tự soạn để mở rộng ngân hàng)
% lop: 8
% kc_du_kien: 8DS-HDT-04, 8DS-HDT-01, 8DS-HDT-05
% do_kho_du_kien: 3
Cho hai số $a$, $b$ thoả mãn $a - b = 2$ và $a^2 + ab + b^2 = 7$. Không tìm $a$, $b$, hãy tính:

a) $a^3 - b^3$;

b) $ab$;

c) $a^2 + b^2$.

\begin{solution}
a) Áp dụng hằng đẳng thức hiệu hai lập phương: $a^3 - b^3 = (a - b)(a^2 + ab + b^2) = 2 \cdot 7 = 14$.

b) Theo hằng đẳng thức bình phương của một hiệu, $(a - b)^2 = a^2 - 2ab + b^2$. Viết $a^2 - 2ab + b^2 = (a^2 + ab + b^2) - 3ab = 7 - 3ab$.

b) Vì $(a - b)^2 = 2^2 = 4$ nên $7 - 3ab = 4$, suy ra $3ab = 3$ và $ab = 1$.

c) Từ giả thiết $a^2 + ab + b^2 = 7$ và kết quả câu b): $a^2 + b^2 = 7 - ab = 7 - 1 = 6$.

Đáp số: a) $\boxed{14}$; b) $\boxed{1}$; c) $\boxed{6}$
\end{solution}
\end{problem}

% ---------------------------------------------------------------- chứng minh tam thức luôn dương, luôn âm
\begin{problem}[8DS-HDT-042]
% nguon: dummy (tự soạn để mở rộng ngân hàng)
% lop: 8
% kc_du_kien: 8DS-HDT-06, 8DS-HDT-01
% do_kho_du_kien: 3
Chứng minh rằng với mọi số thực $x$:

a) $x^2 - 5x + 7 > 0$;

b) $-x^2 + 3x - 4 < 0$.

\begin{solution}
a) Tách $-5x = -2 \cdot x \cdot \dfrac{5}{2}$ và $7 = \dfrac{25}{4} + \dfrac{3}{4}$ để làm xuất hiện bình phương của một hiệu: $x^2 - 5x + 7 = \left(x^2 - 2 \cdot x \cdot \dfrac{5}{2} + \dfrac{25}{4}\right) + \dfrac{3}{4} = \left(x - \dfrac{5}{2}\right)^2 + \dfrac{3}{4}$.

a) Vì $\left(x - \dfrac{5}{2}\right)^2 \geq 0$ nên $x^2 - 5x + 7 \geq \dfrac{3}{4} > 0$ với mọi $x$.

b) Đặt dấu trừ ra ngoài các hạng tử chứa $x$ rồi thêm và bớt $\dfrac{9}{4}$: $-x^2 + 3x - 4 = -\left(x^2 - 2 \cdot x \cdot \dfrac{3}{2} + \dfrac{9}{4}\right) + \dfrac{9}{4} - 4 = -\left(x - \dfrac{3}{2}\right)^2 - \dfrac{7}{4}$.

b) Vì $\left(x - \dfrac{3}{2}\right)^2 \geq 0$ nên $-\left(x - \dfrac{3}{2}\right)^2 \leq 0$, do đó $-x^2 + 3x - 4 \leq -\dfrac{7}{4} < 0$ với mọi $x$.

Đáp số: a) $\boxed{x^2 - 5x + 7 = \left(x - \dfrac{5}{2}\right)^2 + \dfrac{3}{4} > 0}$; b) $\boxed{-x^2 + 3x - 4 = -\left(x - \dfrac{3}{2}\right)^2 - \dfrac{7}{4} < 0}$
\end{solution}
\end{problem}

% ---------------------------------------------------------------- chứng minh chia hết bằng hiệu hai bình phương
\begin{problem}[8DS-HDT-043]
% nguon: dummy (tự soạn để mở rộng ngân hàng)
% lop: 8
% kc_du_kien: 8DS-HDT-02, 8DS-HDT-06
% do_kho_du_kien: 3
Chứng minh rằng với mọi số nguyên $n$:

a) $(n + 7)^2 - (n - 5)^2$ chia hết cho $24$;

b) $(5n + 3)^2 - (5n - 3)^2$ chia hết cho $60$.

\begin{solution}
a) Áp dụng hằng đẳng thức hiệu hai bình phương $A^2 - B^2 = (A - B)(A + B)$ với $A = n + 7$, $B = n - 5$: $(n + 7)^2 - (n - 5)^2 = \left[(n + 7) - (n - 5)\right]\left[(n + 7) + (n - 5)\right]$.

a) Bỏ ngoặc và thu gọn trong mỗi thừa số: $(n + 7) - (n - 5) = 12$ và $(n + 7) + (n - 5) = 2n + 2$. Do đó biểu thức bằng $12(2n + 2) = 24(n + 1)$.

a) Vì $n$ là số nguyên nên $n + 1$ là số nguyên, suy ra $24(n + 1)$ chia hết cho $24$.

b) Áp dụng hằng đẳng thức hiệu hai bình phương với $A = 5n + 3$, $B = 5n - 3$: $(5n + 3)^2 - (5n - 3)^2 = \left[(5n + 3) - (5n - 3)\right]\left[(5n + 3) + (5n - 3)\right] = 6 \cdot 10n = 60n$.

b) Vì $n$ là số nguyên nên $60n$ chia hết cho $60$.

Đáp số: a) $\boxed{(n + 7)^2 - (n - 5)^2 = 24(n + 1) \text{ chia hết cho } 24}$; b) $\boxed{(5n + 3)^2 - (5n - 3)^2 = 60n \text{ chia hết cho } 60}$
\end{solution}
\end{problem}

% ---------------------------------------------------------------- rút gọn rồi tính giá trị với hiệu hai lập phương
\begin{problem}[8DS-HDT-044]
% nguon: dummy (tự soạn để mở rộng ngân hàng)
% lop: 8
% kc_du_kien: 8DS-HDT-03, 8DS-HDT-05
% do_kho_du_kien: 3
Cho biểu thức $M = (x + 3)^3 - (x - 3)^3 - 18x^2 - 2x$.

a) Rút gọn $M$.

b) Tính giá trị của $M$ tại $x = -27$.

\begin{solution}
a) Áp dụng hằng đẳng thức lập phương của một tổng: $(x + 3)^3 = x^3 + 9x^2 + 27x + 27$; áp dụng hằng đẳng thức lập phương của một hiệu: $(x - 3)^3 = x^3 - 9x^2 + 27x - 27$.

a) Thay vào $M$ và bỏ ngoặc (trước ngoặc có dấu trừ thì đổi dấu các hạng tử trong ngoặc): $M = x^3 + 9x^2 + 27x + 27 - x^3 + 9x^2 - 27x + 27 - 18x^2 - 2x$.

a) Thu gọn các hạng tử đồng dạng: $M = (x^3 - x^3) + (9x^2 + 9x^2 - 18x^2) + (27x - 27x - 2x) + (27 + 27) = 54 - 2x$.

b) Thay $x = -27$ vào biểu thức đã rút gọn: $M = 54 - 2 \cdot (-27) = 54 + 54 = 108$.

Đáp số: a) $\boxed{M = 54 - 2x}$; b) $\boxed{M = 108}$
\end{solution}
\end{problem}

% ---------------------------------------------------------------- tích bốn số tự nhiên liên tiếp cộng 1
\begin{problem}[8DS-HDT-045]
% nguon: dummy (tự soạn để mở rộng ngân hàng)
% lop: 8
% kc_du_kien: 8DS-HDT-06, 8DS-HDT-01
% do_kho_du_kien: 4
Chứng minh rằng tích của bốn số tự nhiên liên tiếp cộng với $1$ luôn là bình phương của một số tự nhiên.

\begin{solution}
Gọi bốn số tự nhiên liên tiếp là $n$, $n + 1$, $n + 2$, $n + 3$ với $n$ là số tự nhiên. Ghép thừa số đầu với thừa số cuối và hai thừa số giữa: $n(n + 1)(n + 2)(n + 3) = \left[n(n + 3)\right]\left[(n + 1)(n + 2)\right]$.

Nhân đa thức với đa thức: $n(n + 3) = n^2 + 3n$ và $(n + 1)(n + 2) = n^2 + 3n + 2$.

Đặt $t = n^2 + 3n$, khi đó $n(n + 1)(n + 2)(n + 3) + 1 = t(t + 2) + 1 = t^2 + 2t + 1$.

Áp dụng hằng đẳng thức bình phương của một tổng: $t^2 + 2t + 1 = (t + 1)^2 = (n^2 + 3n + 1)^2$.

Vì $n$ là số tự nhiên nên $n^2 + 3n + 1$ cũng là số tự nhiên. Vậy tích của bốn số tự nhiên liên tiếp cộng với $1$ là bình phương của một số tự nhiên (điều phải chứng minh).

Đáp số: $\boxed{n(n + 1)(n + 2)(n + 3) + 1 = (n^2 + 3n + 1)^2}$
\end{solution}
\end{problem}

% ---------------------------------------------------------------- tổng đan dấu các bình phương liên tiếp
\begin{problem}[8DS-HDT-046]
% nguon: dummy (tự soạn để mở rộng ngân hàng)
% lop: 8
% kc_du_kien: 8DS-HDT-02, 8DS-HDT-05
% do_kho_du_kien: 4
Tính giá trị của $A = 100^2 - 99^2 + 98^2 - 97^2 + \dots + 4^2 - 3^2 + 2^2 - 1^2$.

\begin{solution}
Nhóm hai số hạng liền nhau thành từng cặp: $A = (100^2 - 99^2) + (98^2 - 97^2) + \dots + (4^2 - 3^2) + (2^2 - 1^2)$. Mỗi cặp có dạng $k^2 - (k - 1)^2$ với $k = 100, 98, \dots, 4, 2$.

Áp dụng hằng đẳng thức hiệu hai bình phương: $k^2 - (k - 1)^2 = \left[k - (k - 1)\right]\left[k + (k - 1)\right] = 1 \cdot (2k - 1) = k + (k - 1)$.

Do đó $A = (100 + 99) + (98 + 97) + \dots + (4 + 3) + (2 + 1)$, tức là $A$ bằng tổng của tất cả các số nguyên từ $1$ đến $100$: $A = 1 + 2 + 3 + \dots + 100$.

Tổng của $100$ số hạng đầu của dãy số nguyên dương liên tiếp là $\dfrac{100 \cdot (100 + 1)}{2} = 5\,050$.

Đáp số: $\boxed{A = 5\,050}$
\end{solution}
\end{problem}

% ---------------------------------------------------------------- biết x - y và x^3 - y^3, tìm xy và x, y
\begin{problem}[8DS-HDT-047]
% nguon: dummy (tự soạn để mở rộng ngân hàng)
% lop: 8
% kc_du_kien: 8DS-HDT-03, 8DS-HDT-05, 8DS-HDT-01, 8DS-PTNT-05
% do_kho_du_kien: 4
Cho hai số $x$, $y$ thoả mãn $x - y = 3$ và $x^3 - y^3 = 63$.

a) Tính $xy$.

b) Tính $x^2 + y^2$.

c) Tìm $x$ và $y$.

\begin{solution}
a) Từ hằng đẳng thức lập phương của một hiệu: $(x - y)^3 = x^3 - 3x^2y + 3xy^2 - y^3 = x^3 - y^3 - 3xy(x - y)$, suy ra $x^3 - y^3 = (x - y)^3 + 3xy(x - y)$.

a) Thay $x - y = 3$ và $x^3 - y^3 = 63$: $63 = 3^3 + 3xy \cdot 3 = 27 + 9xy$, suy ra $9xy = 36$ và $xy = 4$.

b) Áp dụng hằng đẳng thức bình phương của một hiệu: $x^2 + y^2 = (x - y)^2 + 2xy = 3^2 + 2 \cdot 4 = 17$.

c) Từ $x - y = 3$ suy ra $x = y + 3$. Thay vào $xy = 4$: $y(y + 3) = 4$, tức là $y^2 + 3y - 4 = 0$.

c) Tách $3y = 4y - y$ rồi nhóm hạng tử: $y^2 + 4y - y - 4 = y(y + 4) - (y + 4) = (y + 4)(y - 1) = 0$. Do đó $y = -4$ hoặc $y = 1$, suy ra $x = y + 3 = -1$ hoặc $x = 4$.

c) Thử lại: với $(x; y) = (4; 1)$ có $4 - 1 = 3$ và $4^3 - 1^3 = 63$; với $(x; y) = (-1; -4)$ có $-1 + 4 = 3$ và $(-1)^3 - (-4)^3 = -1 + 64 = 63$ (đều đúng).

Đáp số: a) $\boxed{xy = 4}$; b) $\boxed{x^2 + y^2 = 17}$; c) $\boxed{(x; y) \in \{(4; 1), (-1; -4)\}}$
\end{solution}
\end{problem}

% ---------------------------------------------------------------- cực trị và bất đẳng thức khi biết a + b = 1
\begin{problem}[8DS-HDT-048]
% nguon: dummy (tự soạn để mở rộng ngân hàng)
% lop: 8
% kc_du_kien: 8DS-HDT-06, 8DS-HDT-01, 8DS-HDT-04
% do_kho_du_kien: 4
Cho hai số $a$, $b$ thoả mãn $a + b = 1$.

a) Tìm giá trị nhỏ nhất của $a^2 + b^2$.

b) Chứng minh rằng $a^3 + b^3 = 1 - 3ab$, rồi suy ra $a^3 + b^3 \geq \dfrac{1}{4}$.

\begin{solution}
a) Từ $a + b = 1$ suy ra $b = 1 - a$, do đó $a^2 + b^2 = a^2 + (1 - a)^2 = 2a^2 - 2a + 1$.

a) Tách để làm xuất hiện bình phương của một hiệu: $2a^2 - 2a + 1 = 2\left(a^2 - a + \dfrac{1}{4}\right) + \dfrac{1}{2} = 2\left(a - \dfrac{1}{2}\right)^2 + \dfrac{1}{2}$.

a) Vì $\left(a - \dfrac{1}{2}\right)^2 \geq 0$ nên $a^2 + b^2 \geq \dfrac{1}{2}$. Dấu bằng xảy ra khi và chỉ khi $a = \dfrac{1}{2}$, khi đó $b = \dfrac{1}{2}$. Thử lại: $\dfrac{1}{4} + \dfrac{1}{4} = \dfrac{1}{2}$ (đúng).

b) Áp dụng hằng đẳng thức tổng hai lập phương: $a^3 + b^3 = (a + b)(a^2 - ab + b^2) = a^2 - ab + b^2$. Theo hằng đẳng thức bình phương của một tổng, $a^2 + b^2 = (a + b)^2 - 2ab = 1 - 2ab$, nên $a^3 + b^3 = 1 - 2ab - ab = 1 - 3ab$.

b) Từ hai hằng đẳng thức bình phương của một tổng và của một hiệu: $(a + b)^2 - (a - b)^2 = 4ab$, suy ra $4ab = 1 - (a - b)^2 \leq 1$ (vì $(a - b)^2 \geq 0$), tức là $ab \leq \dfrac{1}{4}$.

b) Do đó $a^3 + b^3 = 1 - 3ab \geq 1 - \dfrac{3}{4} = \dfrac{1}{4}$ (điều phải chứng minh); dấu bằng xảy ra khi $a = b = \dfrac{1}{2}$.

Đáp số: a) $\boxed{\min (a^2 + b^2) = \dfrac{1}{2} \text{ khi } a = b = \dfrac{1}{2}}$; b) $\boxed{a^3 + b^3 = 1 - 3ab \geq \dfrac{1}{4}}$
\end{solution}
\end{problem}

% ---------------------------------------------------------------- nghiệm nguyên dương của hiệu hai bình phương
\begin{problem}[8DS-HDT-049]
% nguon: dummy (tự soạn để mở rộng ngân hàng)
% lop: 8
% kc_du_kien: 8DS-HDT-02
% do_kho_du_kien: 4
Tìm tất cả các cặp số nguyên dương $(x; y)$ thoả mãn $x^2 - y^2 = 45$.

\begin{solution}
Áp dụng hằng đẳng thức hiệu hai bình phương: $x^2 - y^2 = (x - y)(x + y)$. Phương trình trở thành $(x - y)(x + y) = 45$.

Vì $x$, $y$ nguyên dương nên $x + y$ là số nguyên dương; tích bằng $45 > 0$ nên $x - y$ cũng là số nguyên dương. Ngoài ra $x + y > x - y$ vì $y > 0$.

Ta có $45 = 1 \cdot 45 = 3 \cdot 15 = 5 \cdot 9$, nên với số nhỏ đứng trước có ba trường hợp: $(x - y; x + y) = (1; 45)$, $(3; 15)$ hoặc $(5; 9)$.

Cộng hai phương trình của từng trường hợp: $2x = 46$, $2x = 18$, $2x = 14$, suy ra $x = 23$, $x = 9$, $x = 7$; từ $x + y$ suy ra $y = 22$, $y = 6$, $y = 2$ (đều là số nguyên dương).

Thử lại: $23^2 - 22^2 = 529 - 484 = 45$; $9^2 - 6^2 = 81 - 36 = 45$; $7^2 - 2^2 = 49 - 4 = 45$ (đều đúng).

Đáp số: $\boxed{(x; y) \in \{(23; 22), (9; 6), (7; 2)\}}$
\end{solution}
\end{problem}

% ---------------------------------------------------------------- giá trị nhỏ nhất của tích bốn nhân tử
\begin{problem}[8DS-HDT-050]
% nguon: dummy (tự soạn để mở rộng ngân hàng)
% lop: 8
% kc_du_kien: 8DS-HDT-06, 8DS-HDT-01, 8DS-PTNT-03
% do_kho_du_kien: 5
Tìm giá trị nhỏ nhất của biểu thức $P = x(x + 3)(x + 4)(x + 7)$ và các giá trị của $x$ để $P$ đạt giá trị nhỏ nhất đó.

\begin{solution}
Ghép thừa số đầu với thừa số cuối và hai thừa số giữa (hai cặp có cùng tổng các hạng tử tự do $0 + 7 = 3 + 4$): $P = \left[x(x + 7)\right]\left[(x + 3)(x + 4)\right]$.

Nhân đa thức với đa thức: $x(x + 7) = x^2 + 7x$ và $(x + 3)(x + 4) = x^2 + 7x + 12$. Đặt $t = x^2 + 7x$, khi đó $P = t(t + 12) = t^2 + 12t$.

Thêm và bớt $36$ để làm xuất hiện bình phương của một tổng: $P = (t^2 + 2 \cdot t \cdot 6 + 6^2) - 36 = (t + 6)^2 - 36$.

Vì $(t + 6)^2 \geq 0$ nên $P \geq -36$ với mọi $x$. Dấu bằng xảy ra khi và chỉ khi $t + 6 = 0$, tức là $x^2 + 7x + 6 = 0$.

Nhóm hạng tử: $x^2 + 7x + 6 = x^2 + x + 6x + 6 = x(x + 1) + 6(x + 1) = (x + 1)(x + 6)$. Do đó $x = -1$ hoặc $x = -6$.

Thử lại: tại $x = -1$ có $P = (-1) \cdot 2 \cdot 3 \cdot 6 = -36$; tại $x = -6$ có $P = (-6) \cdot (-3) \cdot (-2) \cdot 1 = -36$. Vậy giá trị nhỏ nhất của $P$ là $-36$.

Đáp số: $\boxed{\min P = -36 \text{ khi } x = -1 \text{ hoặc } x = -6}$
\end{solution}
\end{problem}

% ---------------------------------------------------------------- n^4 + 4 là hợp số
\begin{problem}[8DS-HDT-051]
% nguon: dummy (tự soạn để mở rộng ngân hàng)
% lop: 8
% kc_du_kien: 8DS-HDT-01, 8DS-HDT-02, 8DS-HDT-06
% do_kho_du_kien: 5
Chứng minh rằng với mọi số tự nhiên $n \geq 2$, số $n^4 + 4$ là hợp số.

\begin{solution}
Thêm và bớt $4n^2$ để làm xuất hiện bình phương của một tổng: $n^4 + 4 = (n^4 + 4n^2 + 4) - 4n^2 = (n^2 + 2)^2 - (2n)^2$.

Áp dụng hằng đẳng thức hiệu hai bình phương $A^2 - B^2 = (A - B)(A + B)$ với $A = n^2 + 2$, $B = 2n$: $n^4 + 4 = (n^2 - 2n + 2)(n^2 + 2n + 2)$.

Hai thừa số trên đều là số nguyên dương. Viết thừa số thứ nhất $n^2 - 2n + 2 = (n - 1)^2 + 1$; vì $n \geq 2$ nên $n - 1 \geq 1$, do đó $(n - 1)^2 + 1 \geq 2$.

Thừa số thứ hai lớn hơn thừa số thứ nhất: $(n^2 + 2n + 2) - (n^2 - 2n + 2) = 4n > 0$. Do đó thừa số thứ hai cũng lớn hơn $2$.

Vậy $n^4 + 4$ là tích của hai số tự nhiên đều lớn hơn $1$ nên là hợp số. Kiểm tra: $n = 2$ cho $20 = 2 \cdot 10$; $n = 3$ cho $85 = 5 \cdot 17$.

Đáp số: $\boxed{n^4 + 4 = (n^2 - 2n + 2)(n^2 + 2n + 2) \text{ là hợp số khi } n \geq 2}$
\end{solution}
\end{problem}

\end{document}
